\scanpage{13}\mspage{12}

\uline{Remarque} Dans l'isom can \unsure{$(*)$}, remplaçons $\delta$ par $\lambda\delta$ ($\lambda\in\mathbb Z$). Les deux mbs sont remplacés respectivement par $\underline{\mathrm{Pic}}^{\lambda^{n-1}\delta'}_{B/S}\simeq(\underline{\mathrm{Pic}}^{\delta'}_{B/S})^{(\lambda^{n-1})}$, et par $\underline{\mathrm{Pic}}^{-\lambda\delta}_{A/S}\overset B\times(A,\lambda^{n-1}\widetilde{\delta'})\simeq(\underline{\mathrm{Pic}}^{-\delta}_{A/S}\overset B\times(A,\widetilde{\delta'}))^{(\lambda^n)}$. Donc on trouve, grâce à l'égalité des deux mbs pour $\lambda\delta$, que

\[
(\underline{\mathrm{Pic}}^{\delta'}_{B/S})^{\otimes(\lambda^n-\lambda^{n-1})}\simeq\text{torseur trivial}\quad\text{(canoniqt)}
\]

Faisant $\lambda=-1$, on trouve

\[
(\underline{\mathrm{Pic}}^{\delta'}_{B/S})^{\otimes2}\simeq\text{torseur trivial}\quad\text{(canoniqt)}
\]

% slanted
\margin{Vérifier que la trivialisation est la trivialisation canonique définie pour \uline{toute} section $\delta'$ de $\underline{\mathrm{NS}}_{B/S}$}

\uline{NB} 2 est le ppcm des entiers de la forme $\lambda^n-\lambda^{n-1}$].

\deleted{Lorsque} Comme $(\delta')'=\chi(\delta)^{n-2}\delta$, on en conclut, en échangeant les rôles de $A$ et $B$

\[
(\underline{\mathrm{Pic}}^\delta_{\deleted{\unsure{B}}A/S})^{\otimes(2\chi(\delta)^{n-2})}\simeq\text{torseur trivial}\quad\text{(canon.!)}
\]

% slanted
\margin{Si $n\geq2$ (tout est trivial si $n\leq1$)}

Lorsque $\chi(\delta)$ est pair, on a même

\[
(\underline{\mathrm{Pic}}^\delta_{A/S})^{\otimes(\chi(\delta)^{n-2})}\simeq\text{torseur trivial}\quad\text{(canoniqt)}
\]

pour $\chi(\delta)$ \uline{pair}, comme on voit par les réflexions qui suivent\ldots

% slanted
\margin{Si $n\geq3$, vérifier que cette section est encore la section canonique de $(\underline{\mathrm{Pic}}_{A/S}^{\delta\chi(\delta)^{n-2}/2})^{\otimes2}$. Si $n\geq2$, $\chi(\delta)$ pair implique que $\delta$ divisible par 2, et il faudrait prouver encore que la trivialisation est la trivialisation canonique.}

\deleted{Lorsque $\chi(\delta)=1$, et que $\delta$ est une polarisation, il faudrait prouver que la\deleted{\unsure{ite}} section de $(\underline{\mathrm{Pic}}^\delta_{A/S})^{\otimes2}$ obtenue correspond à l'isom $\underline{\mathrm{Pic}}^\delta\simeq\underline{\mathrm{Pic}}^{-\delta}$ transformant $\Theta$ en $\Theta^{-}$\ldots\unsure{.}}
